\documentclass{letter}
\usepackage[utf8]{inputenc}
\usepackage{geometry}
\geometry{a4paper, margin=1in}

\address{}
\signature{}

\begin{document}

\begin{letter}{Editor}
\opening{\textbf{Subject:} When the centre ignores the math, extremes thrive}

\medskip

Dear Editor,

Your article suggests that figures such as Nigel Farage and Zack Polanski gain influence largely because political extremes reinforce each other.

From the perspective of someone working in infectious-disease economics, that diagnosis seems incomplete.

Epidemiology studies dynamic networks: the same people who spread infections also work, earn wages and pay rent. In such systems exponential processes dominate. Once growth exceeds a threshold, the structure of the system changes.

Economic processes behave similarly. Compound returns generate exponential growth, yet economic debate often assumes that sustained compounding can coexist with limited redistribution.

In a dynamic network that assumption is unstable. When returns compound while redistribution remains weak, inequality compounds too. The extremes are not the cause of the instability; they are its symptom. The cause is the centre’s refusal to confront the mathematics of compounding.

\closing{Yours faithfully,}

\end{letter}
\end{document}
